\documentclass[11pt]{article}

\usepackage[margin=1in]{geometry}
\usepackage[T1]{fontenc}
\usepackage[utf8]{inputenc}
\usepackage{lmodern}
\usepackage{enumitem}
\usepackage{hyperref}

\hypersetup{
  colorlinks=true,
  linkcolor=blue!50!black,
  urlcolor=blue!50!black,
  citecolor=blue!50!black,
  pdftitle={Discovery Family Prompts for p-vs-np}
}

\title{Discovery Family Prompts for \texttt{p-vs-np}}
\author{Neuro-symbolic AI Group\\Idiap Research Institute}
\date{March 18, 2026}

\begin{document}
\maketitle

\begin{abstract}
This report materializes family-specific discovery prompts from the current blocked frontier. Each prompt pack corresponds to a conceptual breakthrough family such as realization, transport, alignment, explicitness, constructive instantiation, or semantic bridging.
\end{abstract}

\section{Summary}
\begin{itemize}[leftmargin=*,itemsep=0.2em]
\item Each breakthrough family now carries a distinct discovery prompt contract.
\item Discovery should branch by conceptual move type rather than by blocked route label alone.
\item The prompt packs below are intended to be reused by discovery, governance, and expert-handoff layers.
\item Synthetic category-theoretic and HoTT reframing families are added when the blocked frontier shows structural-composition or equivalence-coherence pressure across multiple routes.
\item Cross-cutting epistemic patterns are included to cover discovery moves that are not reducible to one bridge family.
\end{itemize}

\section{abstraction-search}
\textbf{Prompt title.} Abstraction Search Prompt

\textbf{Obstruction class.} mcsp-magnification

\textbf{Required move.} Find a sharper abstraction that turns the blocked route into a typed, inspectable theorem target.

\textbf{Abstraction required.} A theorem-family-specific carrier, invariant, or bridge law that adds new mathematical content.

\textbf{Minimal theorem shape.} State a route-specific theorem that would resolve the blocked interface and make the next proof obligation materially sharper.

\subsection*{System Prompt}
You are searching for the minimal conceptual breakthrough that would turn a blocked route into a typed theorem target. Do not restate the blocked lemma. Identify the sharper carrier, invariant, or bridge law that would make the next theorem materially more specific.

\subsection*{User Task}
Restate the obstruction as a missing abstraction problem, propose 2-4 typed bridge candidates, and state the first theorem-pair-specific lemma each candidate would make possible.

\subsection*{Required Questions}
\begin{itemize}[leftmargin=*,itemsep=0.2em]
\item What typed object would add mathematical content beyond the current schematic carrier?
\item What route-specific theorem would become stateable if that object existed?
\item Which candidate clarifies the negative route if it fails?
\end{itemize}
\subsection*{Search Modes}
\begin{itemize}[leftmargin=*,itemsep=0.2em]
\item typed abstraction search
\item bridge split
\item counterfactual sharpening
\end{itemize}
\section{realization}
\textbf{Prompt title.} Carrier Realization Prompt

\textbf{Obstruction class.} carrier-realization

\textbf{Required move.} Turn a typed but schematic carrier into an explicit mathematical object family with admissibility conditions and a forcing theorem.

\textbf{Abstraction required.} An explicit obstruction family, witness family, surrogate promise problem, or quantitative measure that is realized from the source theorem rather than only named by role.

\textbf{Minimal theorem shape.} If source theorem S holds in regime R, then one can construct an explicit carrier C with admissibility conditions A such that C preserves invariant I and supports a route-specific forcing lemma to target T.

\subsection*{System Prompt}
You are converting a typed but schematic carrier into an explicit mathematical object family. Do not rename the carrier. Produce admissibility conditions, a realizability route from the source theorem, and the forcing lemma needed at the target.

\subsection*{User Task}
Construct or sharply constrain an explicit witness family, obstruction family, surrogate promise problem, or quantitative measure that realizes the blocked route from source theorem to target theorem.

\subsection*{Required Questions}
\begin{itemize}[leftmargin=*,itemsep=0.2em]
\item What is the concrete carrier set, map, family, or measure?
\item How is it built from the source theorem rather than postulated?
\item What admissibility conditions prevent it from being only a relabeling?
\end{itemize}
\subsection*{Search Modes}
\begin{itemize}[leftmargin=*,itemsep=0.2em]
\item explicit carrier construction
\item equivalence search between live carrier encodings
\item surrogate problem design with forcing semantics
\end{itemize}
\section{transport}
\textbf{Prompt title.} Transport Principle Prompt

\textbf{Obstruction class.} formula-size-transport

\textbf{Required move.} Find a theorem that preserves the relevant invariant across a cross-model move into formula-size.

\textbf{Abstraction required.} A transport law, compilation theorem, or complexity measure that is stable under the source-to-target model shift.

\textbf{Minimal theorem shape.} For every admissible carrier C in the nonuniform-circuit regime, there is a transformed carrier C' in the formula-size regime preserving invariant I and retaining enough quantitative strength to imply target T.

\subsection*{System Prompt}
You are searching for an invariant-preserving transport theorem across a model shift. Do not strengthen the source theorem unless it changes the transport mechanism. Focus on the cross-model map and the preserved quantitative parameter.

\subsection*{User Task}
State the strongest plausible transport law, compilation theorem, or normal form that could move the current carrier into the target model without losing the invariant needed for the payoff.

\subsection*{Required Questions}
\begin{itemize}[leftmargin=*,itemsep=0.2em]
\item What exact quantity must be preserved across the model shift?
\item What transformed carrier should live in the target model?
\item What is the first obstruction showing why naive transport fails?
\end{itemize}
\subsection*{Search Modes}
\begin{itemize}[leftmargin=*,itemsep=0.2em]
\item cross-model transport principle
\item invariant-preserving compilation theorem
\item representation normal form for formula-size transfer
\end{itemize}
\section{bridge-specificity}
\textbf{Prompt title.} Activation Theorem Prompt

\textbf{Obstruction class.} target-grounding

\textbf{Required move.} Stop treating the target theorem as the open problem and relocate search to the bridge that would activate it.

\textbf{Abstraction required.} No new target theorem; discovery should instead seek the carrier, transfer law, or alignment theorem that makes the target theorem applicable.

\textbf{Minimal theorem shape.} Target theorem T is fixed. The real theorem under search is the local activation theorem A that verifies T's hypotheses from current artifacts.

\subsection*{System Prompt}
The target theorem is already known. Your task is to identify the local activation theorem that would verify its hypotheses from current artifacts. Do not rediscover the target theorem.

\subsection*{User Task}
State the activation theorem under search, name the exact target hypothesis that is still unverified, and identify the local bridge object or transfer law needed to verify it.

\subsection*{Required Questions}
\begin{itemize}[leftmargin=*,itemsep=0.2em]
\item Which target-theorem hypothesis is still open locally?
\item What bridge-local theorem would certify that hypothesis?
\item What novelty remains once the target theorem itself is removed from search?
\end{itemize}
\subsection*{Search Modes}
\begin{itemize}[leftmargin=*,itemsep=0.2em]
\item activation theorem search
\item hypothesis verification search
\item bridge-localization refinement
\end{itemize}
\section{barrier-aware transfer}
\textbf{Prompt title.} Barrier-Aware Transfer Prompt

\textbf{Obstruction class.} local-reduction-barrier

\textbf{Required move.} Design a transfer mechanism that survives known local-reduction barriers, or prove that no such mechanism exists under explicit assumptions.

\textbf{Abstraction required.} A barrier-aware reduction gadget family or an impossibility theorem explaining why the reduction regime cannot be crossed.

\textbf{Minimal theorem shape.} Either construct a local reduction family G satisfying locality bounds L and preserving witness structure W, or prove that any such G would contradict barrier theorem B.

\subsection*{System Prompt}
You are searching under an explicit barrier theorem. Either construct a transfer that survives the barrier's constraints or sharpen the impossibility theorem. Do not use generic hardness language that ignores the reduction regime.

\subsection*{User Task}
Describe a barrier-aware gadget family or impossibility result for the blocked transfer, with locality bounds and preserved witness structure made explicit.

\subsection*{Required Questions}
\begin{itemize}[leftmargin=*,itemsep=0.2em]
\item What exact barrier theorem constrains the route?
\item What locality or regime parameter must any positive construction respect?
\item What negative theorem would be a contribution if positive transfer fails?
\end{itemize}
\subsection*{Search Modes}
\begin{itemize}[leftmargin=*,itemsep=0.2em]
\item barrier-aware gadget design
\item impossibility theorem search
\item reduction regime decomposition
\end{itemize}
\section{alignment}
\textbf{Prompt title.} Alignment Theorem Prompt

\textbf{Obstruction class.} target-alignment

\textbf{Required move.} Identify or prove an equivalence between the current carrier class and the object class required by the target theorem.

\textbf{Abstraction required.} An alignment theorem, correspondence theorem, or mediating normal form linking the local carrier to the target's sparse-problem or distinguisher objects.

\textbf{Minimal theorem shape.} Every admissible carrier C from the local route induces an object D in the target theorem's class, and D preserves the quantitative property needed to activate the target theorem.

\subsection*{System Prompt}
You are aligning a local carrier class with the object class demanded by an existing target theorem. Topical similarity is not enough. Search for equivalence, correspondence, or a mediating normal form.

\subsection*{User Task}
Describe the object-class map from the current carrier to the target theorem's object class, including the quantitative property that must survive the correspondence.

\subsection*{Required Questions}
\begin{itemize}[leftmargin=*,itemsep=0.2em]
\item What is the target theorem's exact object class?
\item How does the local carrier induce an object in that class?
\item Which parameter or sparsity threshold is hardest to preserve?
\end{itemize}
\subsection*{Search Modes}
\begin{itemize}[leftmargin=*,itemsep=0.2em]
\item object-class correspondence search
\item normal-form mediation
\item dual characterization of the target object
\end{itemize}
\section{explicitness}
\textbf{Prompt title.} Explicitness Upgrade Prompt

\textbf{Obstruction class.} pit-upgrade

\textbf{Required move.} Upgrade a semantic or existential algebraic object into an explicit algorithmic derandomization premise.

\textbf{Abstraction required.} An explicit generator, certificate family, or algebraic witness system that is constructive enough for PIT rather than merely semantic.

\textbf{Minimal theorem shape.} If semantic object C exists for the relevant circuit class, then one can uniformly and efficiently extract an explicit structure G that yields deterministic PIT for that class.

\subsection*{System Prompt}
You are upgrading a semantic or existential object into an explicit algorithmic one. Do not accept existence without uniformity, efficiency, or extractability.

\subsection*{User Task}
Search for an explicit generator, certificate family, extractor, or uniform construction rule that turns the blocked semantic object into a derandomization-ready or algorithmically inspectable premise.

\subsection*{Required Questions}
\begin{itemize}[leftmargin=*,itemsep=0.2em]
\item What part of the current object is only existential or semantic?
\item What explicit structure would replace that existential step?
\item What efficiency class is required for the upgraded object to matter?
\end{itemize}
\subsection*{Search Modes}
\begin{itemize}[leftmargin=*,itemsep=0.2em]
\item explicit extractor or generator search
\item certificate map construction
\item semantic-to-algorithmic upgrade theorem
\end{itemize}
\section{constructive instantiation}
\textbf{Prompt title.} Constructive Instantiation Prompt

\textbf{Obstruction class.} schema-to-examples

\textbf{Required move.} Turn an abstract schema into explicit natural example families or prove that the schema cannot instantiate naturally in the current way.

\textbf{Abstraction required.} A construction template, example-generating reduction family, or a canonicalization principle for natural examples.

\textbf{Minimal theorem shape.} From abstract schema S, construct an explicit family of natural instances E\_n satisfying properties P and realizing target theorem T.

\subsection*{System Prompt}
You are turning an abstract schema into explicit mathematical examples. Do not restate the schema or list examples without a construction route.

\subsection*{User Task}
Search for a construction template, reduction family, or canonicalization principle that produces explicit natural instances from the abstract schema.

\subsection*{Required Questions}
\begin{itemize}[leftmargin=*,itemsep=0.2em]
\item What is the first explicit instance family that should be produced?
\item What reduction or construction template yields that family?
\item If the construction fails, what does that say about the schema's naturality?
\end{itemize}
\subsection*{Search Modes}
\begin{itemize}[leftmargin=*,itemsep=0.2em]
\item canonical example construction
\item template instantiation
\item example-preserving reduction search
\end{itemize}
\section{semantic bridge}
\textbf{Prompt title.} Semantic Bridge Prompt

\textbf{Obstruction class.} characterization-bridge

\textbf{Required move.} Bridge a local definitional object to a global characterization theorem.

\textbf{Abstraction required.} A semantic correspondence, duality, witness extraction principle, or completeness notion that turns the local definition into the target biconditional.

\textbf{Minimal theorem shape.} Local definitional object D has property P if and only if global characterization condition Q holds, with both directions proved at the right level of abstraction.

\subsection*{System Prompt}
You are bridging a local definitional object to a global characterization theorem. Syntactic normalization is insufficient. Search for the missing semantic equivalence, duality, or witness-extraction principle.

\subsection*{User Task}
Decompose the biconditional or characterization gap into the minimal semantic implications that would connect the local definition to the global theorem.

\subsection*{Required Questions}
\begin{itemize}[leftmargin=*,itemsep=0.2em]
\item What semantic content is missing from the current definition-level object?
\item Which direction of the characterization is harder and why?
\item What duality or witness principle could mediate the bridge?
\end{itemize}
\subsection*{Search Modes}
\begin{itemize}[leftmargin=*,itemsep=0.2em]
\item semantic correspondence theorem
\item duality or normal form search
\item two-way implication decomposition
\end{itemize}
\section{categorical-reframing}
\textbf{Prompt title.} Categorical Reframing Prompt

\textbf{Obstruction class.} structural-composition-gap

\textbf{Required move.} Recast the blocked frontier as a diagrammatic composition problem: identify objects, morphisms, mediating objects, and the exact commutativity or factorization obligation.

\textbf{Abstraction required.} A structural account of the route in terms of source and target objects, typed morphisms, a mediating universal object, and the factorization or naturality law that would make the route compose.

\textbf{Minimal theorem shape.} There exists a mediating object M and typed morphisms making the route commute; otherwise the failure is certified as no morphism, non-functorial transport, broken naturality, or missing universal property.

\subsection*{System Prompt}
Use category theory as a structural and compositional discovery lens, not as decorative abstraction. Identify the source object, target object, typed morphisms already available, and the diagram that should commute if the route is real. Search for a mediating object with a universal property, test whether the route factors through it, and classify failure as no morphism, non-functorial transport, broken naturality, missing factorization, or failed universal property.

\subsection*{User Task}
Reframe the blocked frontier as a family of diagrammatic composition problems. State candidate objects, morphisms, commutative squares, factorization claims, and universal properties that would make the live routes structurally compose.

\subsection*{Required Questions}
\begin{itemize}[leftmargin=*,itemsep=0.2em]
\item What are the source object, target object, and typed morphisms already present in the route?
\item What diagram should commute if the route is legitimate?
\item What mediating object would be universal for this route?
\item Does the route fail because there is no morphism, because the square does not commute, or because the factorization is missing?
\end{itemize}
\subsection*{Search Modes}
\begin{itemize}[leftmargin=*,itemsep=0.2em]
\item diagram construction
\item commutativity checking
\item factorization search
\item universal-property search
\item adjunction-style dualization
\end{itemize}
\section{equivalence-coherence-reframing}
\textbf{Prompt title.} Equivalence / Coherence Reframing Prompt

\textbf{Obstruction class.} equivalence-coherence-gap

\textbf{Required move.} Recast the blocked frontier as an identity, equivalence, transport, or coherence problem rather than as a missing brute-force bridge.

\textbf{Abstraction required.} A typed account of active formulations as types and terms, an explicit equivalence witness or non-equivalence diagnosis, and the transport law needed to move invariants and theorem hypotheses across that equivalence.

\textbf{Minimal theorem shape.} Either exhibit an equivalence E between the live formulations and prove that the target-relevant invariant transports coherently along E, or prove that no such equivalence witness exists in the current regime.

\subsection*{System Prompt}
Use HoTT as an equivalence and coherence discovery lens, not as a replacement foundation. Type the active concepts and bridge objects, ask whether two formulations are equal, equivalent, or only analogous, search for transportable invariants, and classify failure as no equivalence witness, incoherent transport, or weak equivalence where strict identity is needed.

\subsection*{User Task}
Reframe the blocked frontier as a problem about sameness, equivalence, transport, and coherence. State the candidate equivalence claims, the transport obligations they would induce, and the coherence failures that currently block the route.

\subsection*{Required Questions}
\begin{itemize}[leftmargin=*,itemsep=0.2em]
\item Which active concepts should be treated as types, and which concrete witnesses should be treated as terms?
\item Are the live operationalizations equal, equivalent, incomparable, or only analogous?
\item What invariant or theorem hypothesis must transport along the proposed equivalence?
\item Where does coherence fail when the same route is identified in two different ways?
\end{itemize}
\subsection*{Search Modes}
\begin{itemize}[leftmargin=*,itemsep=0.2em]
\item equivalence witness search
\item transport-law search
\item coherence checking
\item identity-vs-equivalence discrimination
\end{itemize}
\section{invariant-discovery}
\textbf{Prompt title.} Invariant Discovery Prompt

\textbf{Rationale.} Several live routes are still searching for the right preserved quantity or witness measure, so invariant discovery remains a first-class discovery move.

\textbf{Linked families.} abstraction-search, realization, transport, bridge-specificity, barrier-aware transfer, alignment, explicitness, constructive instantiation, semantic bridge

\subsection*{System Prompt}
You are searching for the invariant that the current route is really trying to preserve, expose, or measure. Do not begin from the named bridge object. Begin from the quantity, structure, or monotone property that would survive the route and explain the payoff.

\subsection*{User Task}
Identify 2-4 candidate invariants, state how each is computed or witnessed, and say which one would most sharply reorganize the current blocked route.

\subsection*{Required Questions}
\begin{itemize}[leftmargin=*,itemsep=0.2em]
\item What quantity or structure must survive every legitimate route variant?
\item Which candidate invariant is strong enough to force the payoff but weak enough to be realizable?
\item How would failure of the invariant candidate sharpen the obstruction class?
\end{itemize}
\subsection*{Search Modes}
\begin{itemize}[leftmargin=*,itemsep=0.2em]
\item activation theorem search
\item barrier-aware gadget design
\item bridge split
\item bridge-localization refinement
\item canonical example construction
\item certificate extraction
\item certificate map construction
\item counterfactual sharpening
\item cross-model transport principle
\item dual characterization of the target object
\item duality or normal form search
\item effective construction
\item equivalence search between live carrier encodings
\item example-preserving reduction search
\item explicit carrier construction
\item explicit extractor or generator search
\item generalization by faithful embedding
\item hypothesis verification search
\item impossibility theorem search
\item invariant-preserving compilation theorem
\item normal-form mediation
\item object-class correspondence search
\item parameter-preserving reformulation
\item problem-specific to generic theorem bridge
\item reduction regime decomposition
\item representation normal form for formula-size transfer
\item semantic correspondence theorem
\item semantic-to-algorithmic upgrade theorem
\item surrogate problem design with forcing semantics
\item template instantiation
\item two-way implication decomposition
\item typed abstraction search
\item uniformization
\end{itemize}
\section{regime-localization}
\textbf{Prompt title.} Regime Localization Prompt

\textbf{Rationale.} The frontier repeatedly crosses model or threshold boundaries, so discovery should search for the smallest regime in which the route first becomes true.

\textbf{Linked families.} abstraction-search, realization, transport, bridge-specificity, barrier-aware transfer, alignment, constructive instantiation

\subsection*{System Prompt}
You are searching for the smallest model, threshold, asymptotic regime, or promise-gap setting in which the route might first become true. Do not ask for the full theorem immediately if a more local regime would already be a major contribution.

\subsection*{User Task}
Localize the frontier by proposing restricted regimes, threshold windows, or case splits where the target theorem may first activate, and explain what each localized success would imply for the broader programme.

\subsection*{Required Questions}
\begin{itemize}[leftmargin=*,itemsep=0.2em]
\item Which regime change most plausibly removes the current obstruction without trivializing the result?
\item What is the first nontrivial threshold at which the route may become true?
\item Would a localized theorem already count as an independently valuable subresult?
\end{itemize}
\subsection*{Search Modes}
\begin{itemize}[leftmargin=*,itemsep=0.2em]
\item activation theorem search
\item barrier-aware gadget design
\item bridge split
\item bridge-localization refinement
\item canonical example construction
\item counterfactual sharpening
\item cross-model transport principle
\item equivalence search between live carrier encodings
\item example-preserving reduction search
\item explicit carrier construction
\item generalization by faithful embedding
\item hypothesis verification search
\item impossibility theorem search
\item invariant-preserving compilation theorem
\item parameter-preserving reformulation
\item problem-specific to generic theorem bridge
\item reduction regime decomposition
\item representation normal form for formula-size transfer
\item surrogate problem design with forcing semantics
\item template instantiation
\item typed abstraction search
\end{itemize}
\section{problem-representation-shift}
\textbf{Prompt title.} Problem Representation Shift Prompt

\textbf{Rationale.} Several blocked routes look misrepresented rather than simply false, so representation change is itself a discovery pattern.

\textbf{Linked families.} abstraction-search, transport, alignment, semantic bridge

\subsection*{System Prompt}
You are testing whether the blocked route is misrepresented rather than false. Search for an alternative encoding, normal form, dual object, or reformulation in which the route becomes mathematically transparent.

\subsection*{User Task}
Propose representation changes that preserve the source-side content but expose a better target interface, and state the first theorem each representation would make possible.

\subsection*{Required Questions}
\begin{itemize}[leftmargin=*,itemsep=0.2em]
\item Which part of the current representation hides the real inferential structure?
\item What alternative object language or normal form would expose the missing bridge?
\item What mathematical content must stay invariant under the representation change?
\end{itemize}
\subsection*{Search Modes}
\begin{itemize}[leftmargin=*,itemsep=0.2em]
\item bridge split
\item counterfactual sharpening
\item cross-model transport principle
\item dual characterization of the target object
\item duality or normal form search
\item generalization by faithful embedding
\item invariant-preserving compilation theorem
\item normal-form mediation
\item object-class correspondence search
\item parameter-preserving reformulation
\item problem-specific to generic theorem bridge
\item representation normal form for formula-size transfer
\item semantic correspondence theorem
\item two-way implication decomposition
\item typed abstraction search
\end{itemize}
\section{counterexample-led-concept-formation}
\textbf{Prompt title.} Counterexample-Led Concept Formation Prompt

\textbf{Rationale.} The frontier is heavily blocked, so near-counterexamples and barrier analogs should be used to induce better concepts rather than only to reject claims.

\textbf{Linked families.} abstraction-search, realization, transport, barrier-aware transfer, alignment, semantic bridge

\subsection*{System Prompt}
Use failure as a discovery tool. Search for near-counterexamples, barrier analogs, and route-breaking examples, then ask what concept would separate the surviving cases from the failures.

\subsection*{User Task}
Identify the strongest nearby failure patterns, describe what they invalidate, and derive a refined concept or admissibility condition from those failures.

\subsection*{Required Questions}
\begin{itemize}[leftmargin=*,itemsep=0.2em]
\item What is the nearest counterexample or barrier analog to the current route?
\item Which current concept fails to distinguish the positive and negative cases?
\item What new concept becomes visible only after inspecting the failure pattern?
\end{itemize}
\subsection*{Search Modes}
\begin{itemize}[leftmargin=*,itemsep=0.2em]
\item barrier-aware gadget design
\item bridge split
\item counterfactual sharpening
\item cross-model transport principle
\item dual characterization of the target object
\item duality or normal form search
\item equivalence search between live carrier encodings
\item explicit carrier construction
\item generalization by faithful embedding
\item impossibility theorem search
\item invariant-preserving compilation theorem
\item normal-form mediation
\item object-class correspondence search
\item parameter-preserving reformulation
\item problem-specific to generic theorem bridge
\item reduction regime decomposition
\item representation normal form for formula-size transfer
\item semantic correspondence theorem
\item surrogate problem design with forcing semantics
\item two-way implication decomposition
\item typed abstraction search
\end{itemize}
\section{route-decomposition}
\textbf{Prompt title.} Route Decomposition Prompt

\textbf{Rationale.} High-pressure obstruction classes are carrying multiple inferential roles at once, so route decomposition is needed before more local search.

\textbf{Linked families.} abstraction-search, realization, transport, bridge-specificity

\subsection*{System Prompt}
The current route may be carrying too many inferential roles at once. Decompose it into smaller interfaces, each with its own object, invariant, and theorem target.

\subsection*{User Task}
Split the blocked route into a sequence of subinterfaces, identify the minimal theorem at each interface, and state which subresult would already be valuable if the full route remained open.

\subsection*{Required Questions}
\begin{itemize}[leftmargin=*,itemsep=0.2em]
\item Which inferential roles are currently being conflated in one bridge object?
\item What is the smallest theorem that isolates each role?
\item Which decomposed segment is most likely to produce a standalone contribution?
\end{itemize}
\subsection*{Search Modes}
\begin{itemize}[leftmargin=*,itemsep=0.2em]
\item activation theorem search
\item bridge split
\item bridge-localization refinement
\item counterfactual sharpening
\item cross-model transport principle
\item equivalence search between live carrier encodings
\item explicit carrier construction
\item hypothesis verification search
\item invariant-preserving compilation theorem
\item representation normal form for formula-size transfer
\item surrogate problem design with forcing semantics
\item typed abstraction search
\end{itemize}
\section{hybridization-of-partial-bridges}
\textbf{Prompt title.} Hybrid Bridge Synthesis Prompt

\textbf{Rationale.} Some blocked routes sit at the intersection of multiple breakthrough families, which is evidence that one bridge object may be insufficient and a composed route should be explored.

\textbf{Linked families.} abstraction-search, realization, transport, bridge-specificity, explicitness, constructive instantiation, semantic bridge

\subsection*{System Prompt}
Do not assume one bridge object must do all the work. Search for a hybrid route in which two partial bridges compose, with each bridge carrying a distinct inferential burden.

\subsection*{User Task}
Identify complementary partial bridges, specify their composition order, and state the intermediate theorem that would let them work together.

\subsection*{Required Questions}
\begin{itemize}[leftmargin=*,itemsep=0.2em]
\item Which blocked routes solve different parts of the same theorem pair?
\item What intermediate object or theorem would allow the partial bridges to compose?
\item What failure would show that the hybrid route is illegitimate rather than merely incomplete?
\end{itemize}
\subsection*{Search Modes}
\begin{itemize}[leftmargin=*,itemsep=0.2em]
\item activation theorem search
\item bridge split
\item bridge-localization refinement
\item canonical example construction
\item certificate extraction
\item certificate map construction
\item counterfactual sharpening
\item cross-model transport principle
\item duality or normal form search
\item effective construction
\item equivalence search between live carrier encodings
\item example-preserving reduction search
\item explicit carrier construction
\item explicit extractor or generator search
\item hypothesis verification search
\item invariant-preserving compilation theorem
\item representation normal form for formula-size transfer
\item semantic correspondence theorem
\item semantic-to-algorithmic upgrade theorem
\item surrogate problem design with forcing semantics
\item template instantiation
\item two-way implication decomposition
\item typed abstraction search
\item uniformization
\end{itemize}

\end{document}
